\documentclass[../report.tex]{subfiles}
\begin{document}

Kiểm thử chia hai mức: test vector cho từng module và chương trình chạy trên toàn CPU. Ở mức CPU, kết quả của
Logisim được so với một mô hình tham chiếu viết độc lập với mạch. Mọi kết quả trong chương này được chạy lại ngày
05/11/2026 trên Logisim-evolution 4.1.0, Java 21, Windows 11.

\section{Toolchain}

\subsection{Assembler \file{tools/rvasm.py}}

Assembler dịch toàn bộ RV32I, các lệnh giả thông dụng (\texttt{li}, \texttt{la}, \texttt{mv}, \texttt{nop},
\texttt{j}, \texttt{call}, \texttt{ret}, \texttt{beqz}, \texttt{bgt}, \ldots), các directive \texttt{.text},
\texttt{.data}, \texttt{.word}, \texttt{.half}, \texttt{.byte}, \texttt{.space}, \texttt{.align}, \texttt{.ascii},
\texttt{.asciz}, \texttt{.equ} và \texttt{\%hi()}, \texttt{\%lo()}. Đầu ra là file \file{.text.hex} và
\file{.data.hex} để nạp vào ROM và RAM bằng \emph{Load Image}, kèm file \file{.lst} ghi địa chỉ, mã máy và dòng
nguồn. Tùy chọn \texttt{--circ} ghi thẳng mã máy vào ROM \sig{IMEM} trong file \file{.circ}. Tùy chọn
\texttt{--nohazard} chèn \texttt{nop} cho bản NH theo quy tắc ở mục~\ref{sec:nh}.

Lệnh giả \texttt{li} với hằng 32 bit được tách thành \texttt{lui} và \texttt{addi}. Vì \texttt{addi} mở rộng dấu
immediate 12 bit, khi bit 11 của phần thấp bằng 1 thì phần cao phải cộng thêm 1; \file{t02\_upper.s} kiểm tra
trường hợp hằng \texttt{0x800}.

\subsection{Golden model \file{tools/rvsim.py}}

\file{rvsim.py} là trình mô phỏng ở mức lệnh, viết bằng Python, không dùng gì từ file mạch. Nó thực thi chương
trình, ghi lại 32 thanh ghi, toàn bộ RAM, chuỗi ký tự ra TTY, giá trị GPIO\_OUT, và đếm số lần load-use và số lần
nhảy để tính số chu kỳ của PL theo công thức ở chương~\ref{ch:pipe}. Với SC, số chu kỳ bằng số lệnh.

\subsection{Bộ chạy test \file{tools/run\_tests.py}}

Với mỗi chương trình, \file{run\_tests.py} làm bốn bước:

\begin{enumerate}
  \item Assemble ở chế độ dump: lệnh EBREAK cuối chương trình được thay bằng một đoạn mã lưu x1 đến x31 vào vùng
        \texttt{0x704}--\texttt{0x77C} rồi mới dừng. Nhờ vậy giá trị thanh ghi đọc được từ file RAM xuất ra.
  \item Ghi mã máy vào ROM của một bản sao file \file{.circ}.
  \item Chạy Logisim không giao diện (\texttt{-t table,halt,tty}), nạp RAM bằng \texttt{-l}, xuất RAM bằng
        \texttt{--save}. Nếu chương trình có dòng \texttt{\# STDIN:}, runner gõ chuỗi đó vào Keyboard.
  \item So với golden model: 31 thanh ghi, toàn bộ 16 KiB RAM, chuỗi TTY, GPIO\_OUT, và hai bộ đếm \sig{cycles},
        \sig{instret}. Riêng chương trình có nhập UART thì không so số chu kỳ, vì thời điểm ký tự đến phụ thuộc
        Logisim.
\end{enumerate}

Vì so cả số chu kỳ, bộ test bắt được lỗi về thời gian mà kết quả cuối vẫn đúng, ví dụ stall thừa một chu kỳ hoặc
flush thiếu.

\section{Test vector từng module}\label{sec:vec}

Mỗi file trong \file{tests/vectors/} có dòng tiêu đề là tên chân và mỗi dòng sau là một bộ giá trị vào cùng giá trị
ra mong đợi. Giá trị mong đợi được sinh từ mô hình Python trong \file{generator/refmodels.py}. Hai module tuần tự
\sig{RegisterFile} và \sig{IOController} có cột \sig{CLK} trong vector để tạo cạnh lên. Bảng~\ref{tab:vec} là kết
quả chạy bằng lệnh \texttt{java -jar logisim-evolution-4.1.0-all.jar -w <Module> <vector> <file.circ>}.

\begin{table}[H]
\centering\small
\caption{Kết quả test vector từng module (số vector đạt / tổng số)}\label{tab:vec}
\begin{tabularx}{\textwidth}{@{}l>{\raggedright\arraybackslash}Xrr@{}}
\toprule
Module & Nội dung vector & SC & PL \\
\midrule
\sig{ALU} & ngẫu nhiên và giá trị biên \texttt{0x7FFFFFFF}, \texttt{0x80000000}, $-1$ & 250/250 & 250/250 \\
\sig{ALUControl} & vét cạn 64 tổ hợp & 64/64 & 64/64 \\
\sig{ImmGen} & lệnh ngẫu nhiên với mọi \sig{ImmSel} & 300/300 & 300/300 \\
\sig{Control} & lệnh ngẫu nhiên, gồm lệnh không hợp lệ và lệnh 0 & 400/400 & 400/400 \\
\sig{RegisterFile} & tuần tự, 200 bước ghi và đọc, mỗi bước 2 dòng clock & 400/400 & 400/400 \\
\sig{BranchUnit} & funct3 ngẫu nhiên, 20\% cặp toán hạng bằng nhau & 300/300 & 300/300 \\
\sig{LoadUnit} & funct3 và offset 0--3 ngẫu nhiên & 300/300 & 300/300 \\
\sig{StoreUnit} & funct3 và offset 0--3 ngẫu nhiên & 300/300 & 300/300 \\
\sig{AddrDecode} & địa chỉ RAM và I/O & 200/200 & 200/200 \\
\sig{IOController} & tuần tự, 200 bước đọc ghi các thanh ghi I/O & 400/400 & 400/400 \\
\sig{ForwardUnit} & số hiệu thanh ghi dồn vào x0, x1, x2, x5 để hay trùng & -- & 400/400 \\
\sig{HazardUnit} & như trên, thêm \sig{Redirect} và \sig{D\_Halt} & -- & 400/400 \\
\sig{IDBypass} & như trên, dữ liệu 32 bit ngẫu nhiên & -- & 300/300 \\
\midrule
Tổng & & 2914/2914 & 4014/4014 \\
\bottomrule
\end{tabularx}
\end{table}

Ba khối \sig{ForwardUnit}, \sig{HazardUnit}, \sig{IDBypass} chỉ có trong file pipeline nên cột SC để trống.

\section{Chương trình test toàn CPU}

\begin{table}[H]
\centering\small
\caption{Các chương trình test trong \file{programs/}}\label{tab:progs}
\begin{tabularx}{\textwidth}{@{}lX@{}}
\toprule
Chương trình & Nội dung \\
\midrule
\file{demo.s} & Chương trình nạp sẵn trong ROM: in \texttt{Hello RISC-V!}, tính tổng 12 số Fibonacci (\texttt{0xE8}), hiện lên LED 7 đoạn và bật 8 LED \\
\file{t01\_alu.s} & Mọi lệnh R và I số học, gồm trường hợp biên và ADDI với immediate âm \\
\file{t02\_upper.s} & LUI, AUIPC và hằng lớn \\
\file{t03\_branch.s} & Sáu lệnh rẽ nhánh, cả nhánh nhảy và không nhảy \\
\file{t04\_jump.s} & JAL, JALR, gọi và trả về hàm, bảng nhảy \\
\file{t05\_mem.s} & LW, SW và các lệnh load, store byte, halfword \\
\file{t06\_hazard.s} & Forwarding khoảng cách 1, 2, 3; load-use; ghi x0; nhảy ngược; JALR dùng thanh ghi vừa load \\
\file{t07\_fib.s} & 20 số Fibonacci đầu tiên ghi vào bộ nhớ \\
\file{t08\_sort.s} & Sắp xếp nổi bọt 10 số có dấu trong \texttt{.data} \\
\file{t09\_recursion.s} & Giai thừa và Fibonacci đệ quy dùng stack \\
\file{t10\_hello.s} & In chuỗi ra TTY \\
\file{t11\_io.s} & Đọc ghi GPIO, nhận ký tự từ Keyboard, đổi sang chữ hoa và in ra TTY \\
\bottomrule
\end{tabularx}
\end{table}

Ngoài 12 chương trình trên, \file{tools/fuzz.py} sinh 25 chương trình ngẫu nhiên (seed 0 đến 24), mỗi chương trình
khoảng 120 lệnh trộn số học, load, store và rẽ nhánh với nhiều phụ thuộc gần nhau để kích hoạt forwarding và stall.

\section{Kết quả}\label{sec:testres}

\begin{table}[H]
\centering\small
\caption{Kết quả chạy \file{run\_tests.py}}\label{tab:testres}
\begin{tabular}{@{}llrr@{}}
\toprule
File mạch & Chế độ assemble & 12 chương trình & 25 chương trình fuzz \\
\midrule
SC & bình thường & 12/12 PASS & 25/25 PASS \\
PL & bình thường & 12/12 PASS & 25/25 PASS \\
NH & bình thường & 0/12 PASS & không chạy \\
NH & \texttt{--nohazard} & 12/12 PASS & 25/25 PASS \\
\bottomrule
\end{tabular}
\end{table}

Ở mọi lần PASS, số chu kỳ Logisim đếm được khớp đúng với golden model, trừ \file{t11\_io.s} là chương trình không so
số chu kỳ. Bản NH chạy chương trình bình thường thì sai ngay từ \file{demo.s}: thanh ghi \texttt{sp} ra
\texttt{0xFFFFFFF0} thay vì \texttt{0x00003FF0}, vì \texttt{addi sp, sp, -16} đọc \texttt{sp} trước khi
\texttt{lui sp, 4} kịp ghi. Kết quả này cho thấy bộ test phát hiện được việc thiếu forwarding và stall.

\end{document}
